\section{Theoretical Framework}

\subsection{The Bio-Gravitational Number: A Stability Threshold}

The central question is: under what physical conditions can an organism maintain its shape against gravity? We derive a dimensionless parameter, the \textbf{Bio-Gravitational Number}:
\begin{equation}
\mathcal{B}_g = \frac{EI}{M g L^2},
\label{eq:bg_number}
\end{equation}
where $EI$ is spinal flexural stiffness, $M$ is body mass, $g$ is gravitational acceleration, and $L$ is spine length. This ratio compares passive structural resistance to gravitational bending load.

Cross-species analysis reveals a critical threshold at $\mathcal{B}_g \approx 0.1$. Species above this value (most quadrupeds) can resist gravity passively through bone and disc stiffness alone. Humans occupy a unique position: adult $\mathcal{B}_g \approx 0.01$, placing us deep within the ``Allometric Trap'' where passive resistance fails and the normal spinal S-curve (cervical lordosis, thoracic kyphosis, lumbar lordosis) requires continuous metabolic maintenance via cellular strain sensors (PIEZO1/2, primary cilia) and proprioceptive feedback.

\subsection{The Energy Deficit Window}

During adolescent growth, two unfavorable scaling relationships collide. The metabolic cost of maintaining the S-curve against gravity scales approximately as $L^3$ (elastic strain energy $U \sim EI\kappa^2 L$ with $I \sim L^4$, $\kappa \sim L^{-1}$). Nutrient supply to avascular intervertebral discs, governed by diffusion and convective transport through endplate surface area, scales as $L^{2}$ (surface-area limited)~\cite{urban2004nutrition, sato2025convective}. This widening deficit requires that paraspinal metabolic supply be \textbf{surface- (perfusion-) limited}: were supply volume-limited ($\sim L^3$), demand and supply would co-scale and no deficit window would open.

The deficit ratio $R(L)=C\,L^{\,n-2}$ with $n=3$ increases monotonically with length; its absolute crossover depends on the composite prefactor $C$ (muscle efficiency, $EI$, perfusion density), which we do not calibrate here. The \textbf{prefactor-free} prediction is relative deficit growth across the growth spurt: $R(0.45)/R(0.35)=(0.45/0.35)^{1}\approx 1.29$, i.e.\ a ${\sim}29\%$ rise in the demand-to-supply ratio over the peak-height-velocity interval. We therefore present the growth-spurt window as the interval of fastest deficit \emph{increase}, not as a single calibrated crossover length. This age-length correspondence matches the well-known clinical peak incidence of Adolescent Idiopathic Scoliosis without parameter fitting to clinical data.

\subsection{The Recovery Ratchet: Restoration versus Growth-Remodeling}

We interpret AIS onset not as a postural-balance failure but as a \emph{geometry-recovery} failure. Growth continually perturbs the spine from its reference shape; the active system must restore that shape faster than growth cements the deviation. We split the standing curvature deviation into a recoverable elastic component $\kappa_e$ (the part that reduces supine or on bending films) and a permanently remodeled component $\kappa_p$. Two rates compete: the active geometric-restoration rate $k_r$ and the growth-remodeling rate $k_g(t)$, by which a sustained deviation is cemented through Hueter--Volkmann growth-plate modulation~\cite{stokes2006hueter, villemure2009growth}. For a perturbation that recoils at rate $k_r$ and is cemented at rate $k_g$, the fraction converted to permanent set is $k_g/(k_r+k_g)$, giving
\begin{equation}
\frac{d\kappa_p}{dt} = k_g(t)\,\kappa_e\,\frac{k_g(t)}{k_r + k_g(t)}, \qquad \mathrm{flexibility}(t) = \frac{\kappa_e}{\kappa_e + \kappa_p(t)}.
\label{eq:recovery_ratchet}
\end{equation}
When $k_r \gg k_g$ the spine recoils and little sets (flexibility $\to 1$); when $k_g \gtrsim k_r$---as occurs when $k_g(t)$ peaks at peak height velocity (PHV)---each perturbation is partly cemented before it can be recovered, and the deviation \emph{ratchets} monotonically into a permanent curve (Fig.~\ref{fig:recovery_ratchet}). Because $k_g(t)$ is maximal during the growth spurt, the ratchet predicts the clustering of AIS onset at PHV \emph{without} invoking a postural-balance deficit, consistent with affected adolescents being otherwise neurologically normal and curves being detected incidentally. The Energy Deficit Window acts within this framework by lowering $k_r$---starving the metabolically expensive restoration machinery---and known genetic risk factors (reduced collagen-dependent damping, altered mechanosensation) likewise shift an individual toward the ratchet regime. The model yields a falsifiable structural signature: \emph{curve flexibility (reducibility) declines as $\kappa_p$ accumulates}, so progressing curves should become progressively less reducible on supine and bending films---a prediction already corroborated by the established clinical association between low curve flexibility and progression~\cite{cheung2020supine, ohrtnissen2016flexibility, alfraihat2022random}.

\subsection{Information-Elasticity Coupling: HOX Genes to Geometry}

The model formalizes how developmental patterning genes (particularly HOX genes encoding vertebral identity along the cranio-caudal axis) translate into mechanical geometry. We introduce an \textbf{Information Field} $I(s)$ representing the spatial concentration of HOX-encoded positional information along the spine (coordinate $s$). This field couples to spine mechanics through three channels (see Supplementary Theory for full derivation):

\begin{enumerate}
\item \textbf{Rest curvature modulation}: Information gradients specify target curvatures (lordotic at cervical/lumbar, kyphotic at thoracic).
\item \textbf{Stiffness modulation}: Information modulates local bending and torsional stiffness via extracellular matrix composition.
\item \textbf{Active moment generation}: Information gradients drive active muscular moments maintaining the S-curve.
\end{enumerate}

When the energy deficit exceeds a critical threshold, this active control system can no longer suppress small left-right asymmetries (inevitable developmental noise, typically $\sim 1$--5\%). Under the recovery ratchet above, such asymmetries---benign during childhood because they are readily recovered---are instead cemented rather than recovered once $k_g \gtrsim k_r$, ratcheting into macroscopic scoliotic deformities (Cobb angles $> 15^{\circ}$) during the deficit window.

\subsection{Why Lateral (Coronal Plane) Buckling?}

The spine is loaded primarily in the sagittal plane (forward-backward) by gravity. Why does deformity emerge laterally? The model predicts a \textbf{directional-stiffness anisotropy}: developmental patterning gradients (encoding the S-curve) are aligned with gravitational stress in the sagittal plane, providing maximal effective stiffness in that direction. Small lateral asymmetries create orthogonality between information and stress vectors, reducing effective lateral stiffness by up to 80\%. This creates a localized ``soft spot'' at the thoracolumbar junction (T8--T10, corresponding to steepest information gradient), directing instability into the coronal plane and explaining the anatomic localization of right thoracic AIS.

\textit{Full mathematical derivations, stability analysis, and parameter sensitivity studies are provided in Supplementary Theory Sections S1--S4.}
